\begin{table}[htbp]
\centering
\small
\caption{Oráculo sobre E3, E1 e E2, Passos 1 a 3: recuperação de uma métrica conhecida pelos dois estimadores.}
\label{tab:bloco1_oraculo_passos_1_a_3}
\begin{tabular}{lllllrr}
\toprule
Problema & W verdadeiro & Pontos & Algoritmo & W recuperado (norm. ΣW) & Cosseno & Acurácia (\%) \\
\midrule
Problema 1 & [1,0; 1,0] & E (150 pts) & perceptron & [0,51; 0,49] & 1,000 (0,999–1,000) & 100,0 \\
Problema 1 & [1,0; 1,0] & E (150 pts) & nnls & [0,50; 0,50] & 0,998 (0,996–1,000) & 99,8 (99,3–100,0) \\
\midrule
Problema 2 & [0,3; 1,5] & E (150 pts) & perceptron & [0,17; 0,83] & 1,000 (1,000–1,000) & 100,0 \\
Problema 2 & [0,3; 1,5] & E (150 pts) & nnls & [0,26; 0,74] & 0,991 (0,989–0,993) & 96,7 (95,3–98,7) \\
\midrule
Problema 3 & [1,5; 0,3] & E (150 pts) & perceptron & [0,80; 0,20] & 0,998 (0,997–1,000) & 100,0 \\
Problema 3 & [1,5; 0,3] & E (150 pts) & nnls & [0,73; 0,27] & 0,989 (0,987–0,990) & 98,2 (96,7–99,3) \\
\bottomrule
\multicolumn{7}{p{0.95\linewidth}}{\footnotesize Teste-oráculo (não há LLM): Passo 1, os rótulos vêm da regra do centróide mais próximo sob o W verdadeiro; Passo 2, os dois estimadores tentam recuperá-lo; Passo 3, comparação por acurácia e cosseno. Problema 1 = Oráculo sobre E3 (W = [1; 1]); Problema 2 = Oráculo sobre E1 (W = [0,3; 1,5]) e Problema 3 = Oráculo sobre E2 (W = [1,5; 0,3]), os três sobre os MESMOS pontos, as gaussianas do Problema E (centróides $(-1{,}5; 1)$ e $(1{,}5; -1)$, que diferem nas duas coordenadas; nos centróides do Problema A, com $x_2 = 0$, $w_2$ não seria identificável). W recuperado normalizado (só a direção importa). Acurácia (\%) = acurácia da métrica recuperada contra o rótulo do oráculo. Médias das sementes (mín–máx).} \\
\end{tabular}
\end{table}
